\documentclass{article}

\usepackage{amsmath,amssymb}
\usepackage{xurl}
\usepackage[hidelinks]{hyperref}
\setlength{\emergencystretch}{2em}

% Shared commands for notes in docs/paper-gaps/.

\DeclareUnicodeCharacter{2080}{\ifmmode{}_{0}\else\textsubscript{0}\fi}
\DeclareUnicodeCharacter{2081}{\ifmmode{}_{1}\else\textsubscript{1}\fi}
\DeclareUnicodeCharacter{2082}{\ifmmode{}_{2}\else\textsubscript{2}\fi}
\DeclareUnicodeCharacter{2099}{\ifmmode{}_{n}\else\textsubscript{n}\fi}

\newcommand{\Lean}{\textsc{Lean}}
\ExplSyntaxOn
\NewDocumentCommand{\leanid}{m}{
  \ifmmode
    \textsf{\detokenize{#1}}
  \else
    \group_begin:
    \urlstyle{sf}
    \tl_set:Nn \l_tmpa_tl {#1}
    \tl_replace_all:Nnn \l_tmpa_tl {\_} {_}
    \exp_args:NV \path \l_tmpa_tl
    \group_end:
  \fi
}
\ExplSyntaxOff
\providecommand{\tr}{\operatorname{tr}}
\newcommand{\tnleanIssueBase}{https://github.com/LionSR/TNLean/issues/}
\newcommand{\tnleanPrBase}{https://github.com/LionSR/TNLean/pull/}
\newcommand{\tnleanDocBase}{https://sirui-lu.com/TNLean/docs/}
\newcommand{\ghissue}[1]{\href{\tnleanIssueBase#1}{GitHub issue \##1}}
\newcommand{\ghpr}[1]{\href{\tnleanPrBase#1}{PR \##1}}
\newcommand{\doclink}[1]{\href{\tnleanDocBase#1.html}{\path{#1}}}

% Dirac notation and the identity operator, as in blueprint/src/macros/common.tex.
% \providecommand keeps note-local definitions authoritative.
\usepackage{dsfont}
\providecommand{\ket}[1]{|#1\rangle}
\providecommand{\bra}[1]{\langle #1|}
\providecommand{\braket}[2]{\langle #1 | #2 \rangle}
\providecommand{\ketbra}[2]{|#1\rangle\!\langle #2|}
\providecommand{\Id}{\mathds{1}}
\providecommand{\one}{\mathds{1}}
\providecommand{\C}{\mathbb{C}}
\providecommand{\MN}[1]{M_{#1}(\C)}
\providecommand{\MD}{M_D(\C)}

% External referencing for paper sources.  The build helper writes sanitized
% auxiliary files under build/paper-aux/ and excludes duplicated source labels.
\usepackage{xr}

% Import a generated paper auxiliary file with a caller-supplied label prefix.
% The candidate paths support builds from docs/paper-gaps/, the repository root,
% and build/paper-aux/ itself.
\newcommand{\importpaperaux}[2]{%
  \IfFileExists{../../build/paper-aux/#2.aux}{%
    \externaldocument[#1]{../../build/paper-aux/#2}%
  }{%
    \IfFileExists{build/paper-aux/#2.aux}{%
      \externaldocument[#1]{build/paper-aux/#2}%
    }{%
      \IfFileExists{#2.aux}{%
        \externaldocument[#1]{#2}%
      }{}%
    }%
  }%
}

% General external reference macros
\newcommand{\paperref}[2]{\ref{#1-#2}}
\newcommand{\papereqref}[2]{(\ref{#1-#2})}

% CPSV16 source (1606.00608 / MPDO-22-12-17-2.tex) external document and shortcuts
\importpaperaux{cpsv16-}{1606.00608/MPDO-22-12-17-2-tnlean}
\newcommand{\cpsvref}[1]{\paperref{cpsv16}{#1}}
\newcommand{\cpsveqref}[1]{\papereqref{cpsv16}{#1}}

% Verdict marker: \gapnote{<kind>}{<status>}. Kind is the policy
% classification (clarification, local-correction, scope-restriction,
% unfaithful, false-source, open-gap); status is open, wip (elimination
% underway), resolved, or historical. Machine-read by the site index;
% prints a verdict line.
\newcommand{\gapnote}[2]{\par\noindent\textbf{Verdict:} #1 (#2).\par}


\title{Wolf Theorem~8.6 --- the Doubly-Substochastic Half of Birkhoff}
\author{}
\date{\today}

\begin{document}
\maketitle

\section{The source statement}

Wolf's Theorem~8.6 (``Birkhoff's theorem'') states two characterizations
in the same theorem environment
\cite[Chapter~8, Theorem~8.6]{Wolf2012Quantum}.  The local source is
\path{Notes/WolfNoteTexSource/ch08_distance_measures.tex}, lines~263--266.

\begin{enumerate}
  \item[\textbf{DS}] The set of $d\times d$ doubly stochastic matrices is the
    convex hull of the $d\times d$ permutation matrices, and its extreme
    points are exactly the permutation matrices.
  \item[\textbf{Sub}] The set of $d\times d$ doubly substochastic matrices is
    the convex hull of the $d\times d$ partial permutation matrices, and its
    extreme points are exactly the partial permutation matrices.
\end{enumerate}

\section{Formalized component}

The doubly-stochastic half (DS) is formalized as two declarations:
\begin{itemize}
  \item \leanid{TNLean.wolf_birkhoff_doublyStochastic_convexHull}
    --- the set of doubly stochastic matrices is the convex hull of the
    permutation matrices, in \doclink{TNLean/Analysis/Birkhoff};
  \item \leanid{TNLean.wolf_birkhoff_extremePoints_doublyStochastic}
    --- the extreme points of the doubly stochastic matrices are exactly the
    permutation matrices, likewise in \doclink{TNLean/Analysis/Birkhoff}.
\end{itemize}
Both are direct reuse of Mathlib's existing Birkhoff theorems
(\leanid{doublyStochastic_eq_convexHull_permMatrix} and
\leanid{extremePoints_doublyStochastic}) without reproof.

\section{Missing component}

The doubly-substochastic half (Sub) has not been formalized.  It requires
the following definitions and results:
\begin{enumerate}
  \item Define the predicate ``doubly substochastic'': a nonnegative square
    matrix whose row sums and column sums are each bounded by~$1$.
  \item Define the partial permutation matrices: square $\{0,1\}$-matrices
    with at most one~$1$ in each row and each column, all other entries~$0$.
  \item Prove that the set of doubly substochastic matrices is the convex hull
    of the partial permutation matrices, and that its extreme points are
    exactly the partial permutation matrices.  This is the substochastic
    analogue of Birkhoff's theorem.
\end{enumerate}
The substochastic polytope is larger (containing the stochastic polytope as a
face) and the extreme-point argument differs from the stochastic case:
permutation matrices are replaced by partial permutation matrices, and the
support of a substochastic matrix may omit rows or columns.

\section{Elimination plan}

\begin{enumerate}
  \item Port the predicate \leanid{doublyStochastic} from Mathlib to a
    ``doubly substochastic'' predicate \leanid{doublySubstochastic}, borrowing
    the nonnegativity, row-sum, and column-sum structure with $\le 1$ bounds
    instead of~$=1$.
  \item Define the set of partial permutation matrices over
    $\operatorname{Fin} d$ and establish that each is doubly substochastic.
  \item Carry out a Birkhoff-style decomposition by completing a
    $d\times d$ doubly substochastic matrix to a $2d\times 2d$ doubly
    stochastic block matrix.  The diagonal off-blocks encode the row and
    column deficits.  A nonnegative lower-right block has row sums equal to
    the original column sums and column sums equal to the original row sums;
    it can be obtained from their normalized outer product when the total
    matrix sum is nonzero, with the zero case treated separately.  Applying
    Birkhoff's theorem to this completion and restricting permutation
    matrices to the original block yields partial permutation matrices.
  \item Prove that the extreme points of the doubly substochastic set are
    exactly the partial permutation matrices, adapting the extreme-point
    argument from the stochastic case.
\end{enumerate}

\section{Verdict}

This is a \emph{scope restriction}.  The doubly-stochastic half of Wolf
Theorem~8.6 (the two Birkhoff characterizations for stochastic matrices) is
fully formalized and linked to the blueprint.  The doubly-substochastic half
(definition, polytope structure, and extreme partial-permutation-matrix
characterization) is not yet formalized.  Tracking: \ghissue{3959}.

\bibliographystyle{alpha}
\bibliography{references}

\end{document}
